\documentclass{article}
\usepackage{graphicx}
\usepackage[margin=1.15in]{geometry}
\usepackage{xcolor}
\usepackage{parskip}
\usepackage{float}
\usepackage{fancyhdr}
\usepackage{array}
\usepackage{booktabs}
\usepackage{enumitem}
\usepackage{amsmath}

\pagestyle{fancy}
\fancyhf{}
\cfoot{\thepage}
\rhead{Lecture 19}
\lhead{MA 214: Applied Statistics (Spring 2026)}
\renewcommand{\headrulewidth}{.4mm}

\newcommand{\blankline}[1]{\underline{\hspace{#1}}}

\usepackage{listings}
\lstset{
  basicstyle=\ttfamily\small,
  columns=fullflexible,
  frame=single,
  framerule=0.4pt,
  breaklines=true,
  showstringspaces=false,
  keywordstyle=\color{blue!60!black},
  commentstyle=\color{green!40!black},
  stringstyle=\color{orange!60!black}
}

\begin{document}

\noindent\textbf{Name:}\blankline{2.25in}\hfill\textbf{BUID:}\blankline{1.55in}

\medskip

\noindent\textbf{Name:}\blankline{2.25in}\hfill\textbf{BUID:}\blankline{1.55in}\newline

\begin{center}
 \Large In-Class Worksheet: MCMC and Diagnostics
\end{center}

\subsection*{Scenario}
We return to the systolic blood pressure (SBP) data from Lecture 18's activity ($n = 20$ patients, $\bar{x} = 126$ mmHg). There we estimated the mean $\mu$ with a Normal--Normal conjugate model, \emph{assuming} the SD was known ($\sigma = 15$). Today we drop that assumption and estimate $\mu$ and $\sigma$ \textbf{jointly}, with independent priors that are \emph{not} conjugate:
\[
\mu \sim N(120,\; 10^2), \qquad \log(\sigma) \sim N(\log 10,\; 0.5^2).
\]
Two unknowns and no conjugacy $\Rightarrow$ no closed-form posterior $\Rightarrow$ we need the computational tools from today.

\medskip
\noindent\textbf{Goals:}
\textbf{(1)} see how grid approximation works (and struggles) in 2D,
\textbf{(2)} read trace plots and understand burn-in,
\textbf{(3)} identify good and bad MCMC chains,
\textbf{(4)} use $\widehat{R}$ to diagnose convergence.

\medskip
\noindent\textbf{Files.} Download \texttt{lecture19\_activity.R} from Blackboard. Its \texttt{Setup} section loads the SBP data and defines three functions you will call: \texttt{log\_posterior}, \texttt{run\_mcmc} (a Metropolis--Hastings sampler), and \texttt{compute\_rhat}. In practice you would use a probabilistic programming system like Stan or PyMC; this is a zero-install stand-in so we can focus on \emph{interpreting} the output.

\subsection*{Setup}
\begin{enumerate}[label={\bfseries (\Alph*)},itemsep=.5em]
	\item \textbf{Determine roles.} You will work in PAIRS. One person is the \textbf{analyst} (runs the code), the other the \textbf{guide} (reads the questions and fills out the worksheet).
	\item \textbf{Open RStudio.} Open \texttt{lecture19\_activity.R} and run the \texttt{Setup} section.
\end{enumerate}

\subsection*{Part 1: Grid Approximation in 2D}
Grid approximation evaluates the posterior on a $K \times K$ grid of $(\mu, \sigma)$ values. You'll see how the resolution $K$ affects the approximation, then record reference summaries from a fine grid.

\begin{enumerate}[label={\bfseries (\Alph*)},itemsep=.6em]
	\item \textbf{Contour plot.} Run \textbf{1A}, which plots the joint posterior on a $101 \times 101$ grid. The result is a \textbf{contour plot}: a bird's-eye view of a 3D surface (posterior density over $\mu$ and $\sigma$). Each contour connects points of equal density, like elevation lines on a map; the innermost contour marks the posterior peak.

	Over what range of $\mu$ does most of the posterior mass lie? What about $\sigma$? \\[3em]

	\newpage
	\item \textbf{Grid convergence.} Run \textbf{1B}, which computes the posterior mean of $\mu$ and $\sigma$ at seven grid sizes ($K = 5, 7, 11, 15, 21, 51, 101$). At what $K$ do the estimates stabilise? Does one parameter need a finer grid than the other --- and if so, why? \\[5em]

	\item \textbf{Reference summaries.} Run \textbf{1C} with $K = 201$. Record the posterior mean and SD for each parameter:

	\vspace{0.3em}
	\begin{center}
	\renewcommand{\arraystretch}{1.4}
	\begin{tabular}{>{\bfseries}p{3cm}cc}
	\toprule
	\textbf{Parameter} & \textbf{Posterior mean} & \textbf{Posterior SD} \\
	\midrule
	$\mu$ & \blankline{4cm} & \blankline{4cm} \\
	$\sigma$ & \blankline{4cm} & \blankline{4cm} \\
	\bottomrule
	\end{tabular}
	\end{center}
	\vspace{0.5em}

	With 2 parameters, a $K = 21$ grid needs $21^2 = 441$ points. How many would you need for 10 parameters? For 20? This is the \textbf{curse of dimensionality} --- and the reason we turn to MCMC. \\[3em]
\end{enumerate}

\subsection*{Part 2: Running MCMC and Burn-in}
A \textbf{trace plot} shows the sampled values of a parameter (y-axis) across iterations (x-axis). When the chain starts far from the region of high posterior probability, the early iterations drift toward it before settling. That initial transient is the \textbf{burn-in}; those draws don't represent the posterior, so we discard them before computing summaries.

\begin{enumerate}[label={\bfseries (\Alph*)},itemsep=.6em]
	\item \textbf{Trace plots from a bad start.} Run \textbf{2A}, which draws $M = 2{,}000$ samples starting from $(\mu_0, \sigma_0) = (140, 25)$ --- far from the high-density region. Examine both trace plots. Where is the burn-in? Roughly how many iterations until the chain reaches the high-density region? \\[4em]

	\item \textbf{Effect of burn-in.} Run \textbf{2B}, which computes the posterior mean of $\mu$ and $\sigma$ after discarding different amounts of burn-in (none, 200, 500, 1000). How do the estimates change as you discard more? Where do they stabilise? \\[4em]

	\item \textbf{MCMC vs.\ grid.} Run \textbf{2C}, which compares MCMC summaries (dropping the first 500 draws) to the grid reference from Part~1. Record:

	\vspace{0.3em}
	\begin{center}
	\renewcommand{\arraystretch}{1.4}
	\begin{tabular}{>{\bfseries}p{2.5cm}cccc}
	\toprule
	& \multicolumn{2}{c}{\textbf{MCMC ($M = 2{,}000$)}} & \multicolumn{2}{c}{\textbf{Grid (Part 1)}} \\
	\cmidrule(lr){2-3}\cmidrule(lr){4-5}
	\textbf{Parameter} & Mean & SD & Mean & SD \\
	\midrule
	$\mu$ & \blankline{2.2cm} & \blankline{2.2cm} & \blankline{2.2cm} & \blankline{2.2cm} \\
	$\sigma$ & \blankline{2.2cm} & \blankline{2.2cm} & \blankline{2.2cm} & \blankline{2.2cm} \\
	\bottomrule
	\end{tabular}
	\end{center}
	\vspace{0.5em}

	How well do the MCMC estimates match the grid? Why might they differ? \\[3em]

	\item \textbf{Running longer.} Run \textbf{2D}, which re-runs the sampler with $M = 10{,}000$ iterations. Record the new summaries (same burn-in):

	\vspace{0.3em}
	\begin{center}
	\renewcommand{\arraystretch}{1.4}
	\begin{tabular}{>{\bfseries}p{2.5cm}cccc}
	\toprule
	& \multicolumn{2}{c}{\textbf{MCMC ($M = 10{,}000$)}} & \multicolumn{2}{c}{\textbf{Grid (Part 1)}} \\
	\cmidrule(lr){2-3}\cmidrule(lr){4-5}
	\textbf{Parameter} & Mean & SD & Mean & SD \\
	\midrule
	$\mu$ & \blankline{2.2cm} & \blankline{2.2cm} & \blankline{2.2cm} & \blankline{2.2cm} \\
	$\sigma$ & \blankline{2.2cm} & \blankline{2.2cm} & \blankline{2.2cm} & \blankline{2.2cm} \\
	\bottomrule
	\end{tabular}
	\end{center}
	\vspace{0.5em}

	Did the longer chain improve the match? Why does chain length matter even after burn-in? (Hint: think about ESS and MCSE.) \\[3em]
\end{enumerate}

\subsection*{Part 3: Good Chains and Bad Chains}

\begin{enumerate}[label={\bfseries (\Alph*)},itemsep=.6em]
	\item \textbf{Step-size comparison.} Run \textbf{3A}, which runs the sampler with three proposal step sizes. For each, record the acceptance rate and describe the $\mu$ trace plot:

	\vspace{0.3em}
	\begin{center}
	\renewcommand{\arraystretch}{1.6}
	\begin{tabular}{>{\bfseries}p{4cm}cp{6.5cm}}
	\toprule
	\textbf{Step size} & \textbf{Acc.\ rate} & \textbf{Describe the trace plot} \\
	\midrule
	$0.05$ [too small] & \blankline{1.6cm} & \\[2.5em]
	$1.5$ [good] & \blankline{1.6cm} & \\[2.5em]
	$15$ [too large] & \blankline{1.6cm} & \\[2.5em]
	\bottomrule
	\end{tabular}
	\end{center}

	\item In your own words, why does the proposal step size matter --- and what goes wrong at each extreme? \\[4em]
\end{enumerate}

\subsection*{Part 4: Multiple Chains and $\widehat{R}$}
One way to check convergence is to run \emph{multiple chains} from different starting values. If they all converge to the same distribution, we gain confidence the sampler works. $\widehat{R}$ formalises this: it compares variation \emph{within} each chain to variation \emph{between} chains. Converged $\Rightarrow \widehat{R}\approx 1$; $\widehat{R} > 1.01$ is a concern, and $> 1.1$ signals badly unreliable estimates.

\begin{enumerate}[label={\bfseries (\Alph*)},itemsep=.6em]
	\item \textbf{Good chains.} Run \textbf{4A}: 4 chains from different starts, good step size. Do the overlaid traces converge to the same region?

	\medskip
	\begin{center}
	\renewcommand{\arraystretch}{1.4}
	\begin{tabular}{>{\bfseries}p{3cm}c}
	\toprule
	\textbf{Parameter} & $\widehat{R}$ \\
	\midrule
	$\mu$ & \blankline{4cm} \\
	$\sigma$ & \blankline{4cm} \\
	\bottomrule
	\end{tabular}
	\end{center}
	\vspace{0.7em}

	\item \textbf{Bad chains.} Run \textbf{4B}: 4 chains from different starts, but a too-small step size. Examine the traces.

	\medskip
	\begin{center}
	\renewcommand{\arraystretch}{1.4}
	\begin{tabular}{>{\bfseries}p{3cm}c}
	\toprule
	\textbf{Parameter} & $\widehat{R}$ \\
	\midrule
	$\mu$ & \blankline{4cm} \\
	$\sigma$ & \blankline{4cm} \\
	\bottomrule
	\end{tabular}
	\end{center}
	\vspace{0.7em}

	\item Compare the two sets of $\widehat{R}$. What went wrong with the bad chains, and how does $\widehat{R}$ detect it? \\[5em]
\end{enumerate}

\newpage

\subsection*{Part 5: Wrap-Up}
\begin{enumerate}[label={\bfseries (\Alph*)},itemsep=.6em]
	\item Suppose all four chains start at the \emph{same} value $(\mu_0, \sigma_0) = (126, 15)$ with a too-small step size. Would $\widehat{R}$ be close to 1, or large? Would the chains actually be reliable? Explain --- and say what this reveals about what $\widehat{R}$ can and cannot catch. \\[5em]

	\item \textbf{Comparing to Lecture 18.} Last time you estimated $\mu$ for this same SBP data with the Normal--Normal conjugate model, \emph{assuming} $\sigma = 15$ was known. Today's MCMC estimates $\mu$ and $\sigma$ jointly. Fill in today's result from your Part~2D run:

	\medskip
	\begin{center}
	\renewcommand{\arraystretch}{1.6}
	\begin{tabular}{>{\bfseries}p{6cm}cc}
	\toprule
	\textbf{Model} & \textbf{Post.\ mean of $\mu$} & \textbf{Post.\ SD of $\mu$} \\
	\midrule
	Lec 18: $\sigma = 15$ assumed known & $125.39$ & $3.18$ \\
	Lec 19: joint $(\mu, \sigma)$ via MCMC & \blankline{3cm} & \blankline{3cm} \\
	\bottomrule
	\end{tabular}
	\end{center}
	\vspace{0.5em}

	The Lecture 18 model \emph{assumed} it knew $\sigma$ exactly. Today's model admits it does not. Which of the two gives a more honest assessment of the uncertainty in $\mu$, and why? (Think about what pretending to know $\sigma$ leaves out.) \\[5em]
\end{enumerate}

\end{document}
